% questions.tex —— 1957 年·全国卷 · 题目内容片段（无答案）
% 由 tools/generate_projects.py 依来源总账自 1957/1957_数学真题_试题.md 生成，勿手改（改源 MD 后重跑）


\begin{gktip}
\textbf{题面说明}：原卷按“一甲、一乙、一丙、一丁、一戊、二、三、四、五”编号。为适配题库统一接口，本文映射为第 1—9 题；原卷答案页中可见的图形在 LaTeX 成品中全部以 TikZ 矢量源重绘。
\end{gktip}

\gkpart{一、小题}

\begin{question}
\gkti{1}（甲）化简
\[
\left(2\frac79\right)^{\frac12}+0.1^{-2}+\left(2\frac{10}{27}\right)^{-\frac23}.
\]

\end{question}
\begin{question}
\gkti{2}（乙）求适合不等式 $x^2+x<2$ 的实数 $x$ 的范围。

\end{question}
\begin{question}
\gkti{3}（丙）求证
\[
\cot22^\circ30'=1+\sqrt2.
\]
（注：$\cot$ 是余切的符号。）

\end{question}
\begin{question}
\gkti{4}（丁）在四面体 $ABCD$ 中，$AC=BD$，$P$、$Q$、$R$、$S$ 依次为棱 $AB$、$BC$、$CD$、$DA$ 的中点。求证 $PQRS$ 为一菱形。

\end{question}
\begin{question}
\gkti{5}（戊）设 $a$、$b$ 为异面二直线，$EF$ 为 $a$、$b$ 的公垂线，$\alpha$ 为过 $EF$ 中点且与 $a$、$b$ 平行的平面，$M$ 为 $a$ 上任一点，$N$ 为 $b$ 上任一点。求证：线段 $MN$ 被平面 $\alpha$ 二等分。

（注：“异面二直线”也叫“不共面二直线”。）

\end{question}
\gkpart{二、大题}

\begin{question}
\gkti{6}解方程组
\[
\begin{cases}
\lg(2x+1)+\lg(y-2)=1,\\
10^{xy}=10^x\cdot10^y.
\end{cases}
\]
（注：$\lg$ 是以 $10$ 为底的对数 $\log_{10}$ 的符号。）

\end{question}
\begin{question}
\gkti{7}若 $\triangle ABC$ 的内切圆半径为 $r$，求证 $BC$ 边上的高 $AD$ 满足
\[
AD=\frac{2r\cos\dfrac B2\cos\dfrac C2}{\sin\dfrac A2}.
\]

\end{question}
\begin{question}
\gkti{8}若 $\triangle ABC$ 为锐角三角形，以 $BC$ 边为直径作圆，并从 $A$ 作此圆的切线 $AD$，与圆切于 $D$，又在 $AB$ 边上取 $AE$ 等于 $AD$，并过 $E$ 作 $AB$ 的垂线与 $AC$ 边的延长线交于 $F$，求证：

（1）$AE:AB=AC:AF$；

（2）$\triangle ABC$ 的面积等于 $\triangle AEF$ 的面积。

\end{question}
\begin{question}
\gkti{9}求证：方程
\[
x^3-(\sqrt2+1)x^2+(\sqrt2-q)x+q=0
\]
的一个根是 $1$。设这个方程的三个根是一个 $\triangle ABC$ 的三个内角的正弦 $\sin A$、$\sin B$、$\sin C$，求 $A$、$B$、$C$ 的度数及 $q$ 的数值。

\end{question}